% Part: second-order-logic
% Chapter: metatheory
% Section: undecidability-and-axiomatizability

\documentclass[../../../include/open-logic-section]{subfiles}

\begin{document}

\olfileid[fr]{sol}{met}{nax}

\olsection{La logique du second ordre n'est pas
  effectivement axiomatisable}

\begin{thm}
\ollabel{thm:sol-undecidable}
La logique du second ordre est indécidable.
\end{thm}

\begin{proof}
Un !!{sentence} du premier ordre est valide en
logique du premier ordre si et seulement s'il
est valide en logique du second ordre. Or la
validité au premier ordre est indécidable.
\end{proof}

\begin{thm}
\ollabel{cor:sol-not-axiomatizable}
Il n'existe aucun système de !!{derivation}
effectif, à la fois correct et complet, pour
la logique du second ordre sous sa sémantique
standard.
\end{thm}

\begin{proof}
Soit $!A$ un !!{sentence} du premier ordre dans
le langage de l'arithmétique. Par la catégoricité
de~$\Th{PA^2}$ établie en
\olref[sol][met][spa]{cor:sol-pa-categorical},
\[
\Sat{N}{!A}\quad\text{si et seulement si}\quad
\Th{PA^2}\Entails !A.
\]
Soit $!P$ la conjonction des neuf axiomes
de~$\Th{PA^2}$. La dernière condition équivaut
à $\Entails !P\lif !A$ : pour toute
!!{structure} arithmétique~$\Struct{M}$,
on a $\Sat{M}{!P\lif !A}$.

Choisissons $z$, $u$, $u'$, $u''$ et $L$ de leurs sortes
respectives, absentes de $!P$ et de~$!A$. Remplaçons dans
$!P\lif !A$ la
constante~$\Obj{0}$ par la !!{variable}
individuelle~$z$, le successeur~$\prime$
par la !!{variable} de fonction unaire~$u$,
$+$ et $\times$ par les !!{variable}s de
fonctions binaires $u'$ et~$u''$, et $<$ par la
!!{variable} de relation binaire~$L$.
Notons $!P'$ et $!A'$ les résultats de ces
substitutions. L'!!{sentence} de logique pure
du second ordre
\[
\lforall[z][\lforall[u][\lforall[u'][\lforall[u''][\lforall[L][(!P' \lif !A')]]]]]
\]
est valide si et seulement si $\Entails !P\lif !A$,
donc si et seulement si $\Sat{N}{!A}$.

Supposons qu'il existe un système de preuve effectif, correct et
complet pour la logique du second ordre. Énumérons ses
démonstrations sans hypothèses et les !!{sentence}s
arithmétiques. Pour chaque paire constituée d'une démonstration
et d'un !!{sentence} arithmétique, comparons la conclusion de
la démonstration à la transformée de l'énoncé. Si elles
coïncident, renvoyons le code de l'énoncé ; sinon, renvoyons
celui de $\Obj{0}=\Obj{0}$. La correction garantit que tous les
codes renvoyés sont ceux d'énoncés vrais. La complétude
garantit que chaque énoncé vrai apparaît ainsi.

Nous obtenons donc une fonction calculable totale dont l'image
est l'ensemble des codes des !!{sentence}s vrais dans~$\Struct{N}$.
La fonction qui associe à chaque indice le code renvoyé est
représentable dans~$\Th{Q}$ par une !!{formula} du premier
ordre~$!B_f(x,y)$. Alors $\lexists[x][!B_f(x,y)]$ définirait
cet ensemble dans~$\Struct{N}$, ce qui contredit le théorème de
Tarski.
\end{proof}

L'effectivité est essentielle à l'argument. Avec une règle
autorisant comme axiome tout énoncé valide, on obtiendrait
un système correct et complet, mais sans énumération
calculable de ses démonstrations.

\end{document}
